% Шаблон тезисов для конференции «Студенческая весна»
% Лаборатория математики ФН1
% Компиляция: pdflatex tezisy-shablon.tex
% Объём — 1–2 страницы, включая список литературы

\documentclass[12pt, a4paper]{article}

\usepackage[T2A]{fontenc}
\usepackage[utf8]{inputenc}
\usepackage[russian]{babel}
\usepackage[margin=20mm]{geometry}
\usepackage{amsmath, amssymb, amsthm}
\usepackage{graphicx}

\pagestyle{empty}
\setlength{\parindent}{1.25cm}

\theoremstyle{plain}
\newtheorem{theorem}{Теорема}
\newtheorem{lemma}[theorem]{Лемма}

\begin{document}

\noindent УДК 519.2

\begin{center}
  {\large\bfseries НАЗВАНИЕ ДОКЛАДА}

  \vspace{0.5em}

  \textit{И.\,О. Фамилия}

  \vspace{0.3em}

  {\small Лаборатория математики ФН1, группа ФН1-11Б\\
  Научный руководитель: должность, И.\,О. Фамилия\\
  \texttt{email@example.ru}}

  \vspace{0.4em}

  {\small Секция 3. Вычислительная математика и моделирование}
\end{center}

\vspace{0.8em}

\noindent\textbf{Ключевые слова:} три--пять терминов через запятую.

\vspace{0.8em}

Во введении двумя-тремя предложениями поясните, к какому классу относится
задача и почему она интересна. Затем приведите строгую постановку.

Рассматривается задача
\begin{equation}\label{eq:main}
  \frac{\partial u}{\partial t} = a^2\frac{\partial^2 u}{\partial x^2} + f(x, t),
  \qquad 0 < x < L,\ t > 0,
\end{equation}
с начальным условием $u(x, 0) = \varphi(x)$ и краевыми условиями
$u(0, t) = u(L, t) = 0$.

\begin{theorem}
  Формулировка полученного результата. Условия перечисляются явно.
\end{theorem}

Кратко опишите идею доказательства либо схему численного метода.
Для вычислительной работы приведите таблицу или график, подтверждающие
заявленный порядок точности.

\begin{table}[h]
  \centering
  \small
  \caption{Погрешность и наблюдаемый порядок точности}
  \begin{tabular}{|c|c|c|}
    \hline
    Шаг сетки $h$ & Погрешность & Порядок \\
    \hline
    $0{,}1$    & $2{,}1\cdot10^{-3}$ & ---      \\
    $0{,}05$   & $5{,}3\cdot10^{-4}$ & $1{,}99$ \\
    $0{,}025$  & $1{,}3\cdot10^{-4}$ & $2{,}03$ \\
    \hline
  \end{tabular}
\end{table}

В заключение сформулируйте вывод: что получено и как это соотносится
с известными результатами.

\vspace{0.8em}

\begin{thebibliography}{9}
  \footnotesize
  \bibitem{samarskii} \emph{Самарский А.\,А.} Теория разностных схем.~---
    М.: Наука, 1989.~--- 616~с.
  \bibitem{article} \emph{Автор А.\,Б.} Название статьи~// Название журнала.~---
    2025.~--- Т.~11, №~3.~--- С.~40--52.
\end{thebibliography}

\end{document}
